\chapterhead{63}{ORR-24}{Basel III: Finalising Post-Crisis Reforms}{1}{2}

\lohead{63}{a}{Explain the elements of the new standardized approach to measure
operational risk capital, including the business indicator, internal loss
multiplier, and loss component, and calculate the operational risk capital
requirement for a bank using this approach.}

\orsub{The standardized approach}

The standardised approach for measuring operational risk represents a combination
of:

\begin{enumerate}
\item[\textcolor{rust}{\sffamily\bfseries (1)}] a financial statement operational
risk exposure proxy, termed the business indicator, or BI;
\item[\textcolor{rust}{\sffamily\bfseries (2)}] the business indicator component
(BIC), which is the product of a regulatory marginal coefficient and BI; and
\item[\textcolor{rust}{\sffamily\bfseries (3)}] operational loss data specific to
an individual bank, reflected in the internal loss multiplier (ILM).
\end{enumerate}

\orsubb{1) Business Indicator (BI)}

Three-year average of a bank's income from interest, leases, and dividends. Banks
are categorised into buckets based on BI size.

\begin{orkeybox}
\[
\mathrm{BI}=\mathrm{ILDC}_{avg}+\mathrm{SC}_{avg}+\mathrm{FC}_{avg}
\]
where
\[
\begin{aligned}
\mathrm{ILDC} &= \text{interest, lease, dividend component}\\
\mathrm{SC}   &= \text{services component}\\
\mathrm{FC}   &= \text{financial component}
\end{aligned}
\]
\end{orkeybox}

\orsubb{2) Business Indicator Component (BIC)}

BI multiplied by a factor based on the bank's BI bucket.

\noindent\begin{minipage}{\textwidth}
\renewcommand{\arraystretch}{1.35}
\begin{tabularx}{\textwidth}{@{}L{0.14\textwidth}YL{0.26\textwidth}@{}}
\rowcolor{navy}\multicolumn{3}{@{}l@{}}{\hspace{4pt}\thd{BI buckets}}\\
\rowcolor{palenavy} \tsub{Bucket} & \tsub{BI range} & \tsub{BIC}\\
{\tabfont \textbf{1}} & {\tabfont \euro0 billion --- \euro1 billion} & {\tabfont $0.12\times \mathrm{BI}$}\\
\rowcolor{paleblue}
{\tabfont \textbf{2}} & {\tabfont \euro1 billion --- \euro30 billion} & {\tabfont $0.15\times \mathrm{BI}$}\\
{\tabfont \textbf{3}} & {\tabfont \euro30 billion --- $+\infty$} & {\tabfont $0.18\times \mathrm{BI}$}\\
\end{tabularx}
\end{minipage}
\orfigcap{The three BI buckets and their marginal coefficients.}

\begin{orexambox}{Example: computing the BIC}
If BI is equal to \euro40 billion, the BIC would be computed as:
\[
(0.12\times1)+\bigl[0.15\times(30-1)\bigr]+\bigl[0.18\times(40-30)\bigr]
=0.12+4.35+1.8=\text{\euro}6.27\text{ billion}
\]
\end{orexambox}

\orsubb{3) Internal Loss Multiplier (ILM)}

Reflects a bank's historical operational losses relative to the industry average.
Higher losses mean a higher multiplier.

\begin{orkeybox}
\[
\text{internal loss multiplier}=\ln\!\left[e^{1}-1+
\left(\frac{\text{loss component}}{\mathrm{BIC}}\right)^{\!0.8}\right]
\]
where
\[
\text{loss component}=15\times\text{average annual operational risk loss over the last 10 years}
\]
\end{orkeybox}

\begin{ornotebox}{Data requirement}
Minimum 5 years of operational loss data. With less than 5 years, no ILM is
possible.
\end{ornotebox}

The loss component serves to reflect the operational loss exposure based on a
bank's internal loss experiences. If a bank's loss experience is greater (less)
than the industry average, its loss component will be above (below) the BI
component, and its standardised approach capital will be above (below) the BI
component.

\orsub{Standardized approach capital requirement}

\begin{center}
\begin{tikzpicture}[font=\fontsize{8.6}{10.6}\selectfont,
  bx/.style={text width=6.8cm,align=center,rounded corners=3pt,inner sep=9pt,
             text=navy,draw=palenavy,line width=0.7pt}]
\node[bx,fill=paleblue] at (-3.7,0)
  {\textbf{BI bucket 1 banks}\\[3pt]
   operational risk capital $=\mathrm{BIC}$};
\node[bx,fill=palerust] at ( 3.7,0)
  {\textbf{BI bucket 2 and 3 banks}\\[3pt]
   operational risk capital $=\mathrm{BIC}\times \mathrm{ILM}$};
\end{tikzpicture}
\end{center}
\orfigcap{Capital depends on which bucket the bank sits in.}

% ---------------------------------------------------------------------
\lohead{63}{b}{Compare the Standardized Measurement Approach (SMA) to earlier
methods of calculating operational risk capital, including the Advanced
Measurement Approaches (AMA).}

\noindent\begin{minipage}{\textwidth}
\renewcommand{\arraystretch}{1.35}
\begin{tabularx}{\textwidth}{@{}L{0.17\textwidth}YY@{}}
\rowcolor{navy} \thd{} & \thd{AMA} & \thd{SMA}\\
{\tabfont \tsub{Introduction}} &
{\tabfont Part of Basel II, 2006} &
{\tabfont Developed to address AMA's shortcomings}\\
\rowcolor{paleblue}
{\tabfont \tsub{Methodology}} &
{\tabfont Allows use of internal models based on the bank's own data} &
{\tabfont Non-model-based; combines financial info with internal loss data}\\
{\tabfont \tsub{Flexibility}} &
{\tabfont High; banks can develop tailored risk measurement models} &
{\tabfont Low; standardised calculation method for all banks}\\
\rowcolor{paleblue}
{\tabfont \tsub{Challenges}} &
{\tabfont Lack of comparability, complex calculations} &
{\tabfont Aims to simplify and standardise risk calculations}\\
{\tabfont \tsub{Objective}} &
{\tabfont Tailor operational risk management to individual bank needs} &
{\tabfont Enhance comparability and simplicity across the banking sector}\\
\end{tabularx}
\end{minipage}
\orfigcap{AMA versus SMA, across five dimensions.}

% ---------------------------------------------------------------------
\lohead{63}{c}{Describe general and specific criteria recommended by the Basel
Committee for the identification, collection and treatment of operational loss
data.}

\orsub{Identification, collection, and treatment of operational loss data}

\orsubb{General criteria}

\begin{lettered}
\item Banks must have established procedures for identifying, collecting, and
handling operational loss data. This includes dates of occurrence, discovery,
accounting, recoveries, and causes of loss. Loss data must be comprehensive and
reviewed independently.
\item Exclusion of credit risk, but include market risk for operational risk
capital.
\item Banks must classify loss data into Basel II's Level 1 supervisory categories
and document allocation to specific event types.
\item Data must capture all significant exposures and activities across all regions
and systems.
\end{lettered}

\orsubb{Specific criteria}

\begin{lettered}
\item Banks must have clear criteria for when a loss event is included in the loss
data.
\item Banks need to identify gross loss amounts, including and excluding insurance
and non-insurance recoveries.
\end{lettered}
